% =====================================================================
% Apartado 1: Presentación del sistema (1 diapositiva, ~1 minuto)
% Responsable y expositor: Apunte Pazmiño Washington German
% Claves del .bib: spotifyAbout2026, spotifyQ32026, vansteen2023ds,
%                  strid2023fleet
% Cifras del 3T 2026: tabla «User & Financial Summary» de la sección
% Executive Summary del informe (p. 3 según su índice).
% Libro de van Steen: cap. 2, «Cloud computing». Cuando tengas la página
% exacta de tu PDF, cambia [cap.~2] por [p.~XX].

% =====================================================================
\section{Presentación del sistema}

\begin{frame}{Spotify: un servicio de audio a escala mundial}
  \begin{itemize}
    \item \textbf{Resuelve:} descubrir y escuchar música, pódcasts y
      audiolibros en un solo lugar \cite{spotifyAbout2026}.
    \item \textbf{Lo usan:} oyentes gratis con anuncios y suscriptores
      que pagan Premium \cite[p.~3]{spotifyQ32026}.
  \end{itemize}

  \medskip
  \centering
  \begin{tabular}{@{}lr@{}}
    \toprule
    \textbf{Escala en el 3.\textsuperscript{er} trimestre de 2026} &
      \textbf{Cifra (crecimiento anual)} \\
    \midrule
    Usuarios activos al mes (MAU)  & 788 millones (+11\,\%) \cite[p.~3]{spotifyQ32026} \\
    Suscriptores Premium           & 306 millones (+9\,\%) \cite[p.~3]{spotifyQ32026} \\
    Usuarios gratis con anuncios   & 499 millones (+12\,\%) \cite[p.~3]{spotifyQ32026} \\
    Mercados donde está disponible & 184 \cite{spotifyAbout2026} \\
    \bottomrule
  \end{tabular}

  \medskip
  \raggedright
  {\footnotesize \textbf{Nube:} el proveedor ofrece su infraestructura, su
  plataforma o sus aplicaciones como servicio (IaaS, PaaS, SaaS)
  \cite[p.~99]{vansteen2023ds}. Spotify trabaja con Google Cloud desde
  2016 y trasladó allí su plataforma con \emph{lift and shift}
  \cite{strid2023fleet}.\par}
\end{frame}
